%% IEEE variant of accounting_state_model.tex -- generated by
%% scripts/make_ieee_tables.py conventions; the wide figure becomes a full-width
%% float so the two-column layout does not shrink it to illegibility.
% Accounting state model -- included by main.tex (% Accounting state model -- included by main.tex (% Accounting state model -- included by main.tex (% Accounting state model -- included by main.tex (\input{accounting_state_model})
% Requires: tikz (with positioning, arrows.meta), amsmath, booktabs.

\section{Accounting State Model}
\label{sec:statemodel}

We model a metered inference request as a state machine so that each mechanism can be
localized to a specific transition, establishing that the taxonomy is
\emph{architectural} rather than a list of unrelated implementation bugs.

\begin{figure}[t]
\centering
\resizebox{\linewidth}{!}{%
\begin{tikzpicture}[
  node distance=8mm and 12mm,
  every node/.style={font=\scriptsize},
  st/.style={draw, rounded corners=2pt, minimum height=6mm, minimum width=17mm, align=center},
  val/.style={st, fill=black!6},
  com/.style={st, fill=black!12},
  bad/.style={st, draw=red!70!black, thick},
  ->/.default=,
  arr/.style={-{Latex[length=1.6mm]}, thin},
  barr/.style={-{Latex[length=1.6mm]}, thin, red!70!black, dashed},
]
\node[st] (created) {CREATED};
\node[st, right=of created] (auth) {AUTHORIZED};
\node[com, right=of auth] (res) {RESERVED};
\node[st, right=of res] (exec) {EXECUTING};
\node[val, right=of exec] (stream) {STREAMING};
\node[val, above right=4mm and 10mm of stream] (comp) {COMPLETED};
\node[bad, below right=4mm and 10mm of stream] (abort) {ABORTED};
\node[st, right=14mm of stream, xshift=26mm] (acc) {ACCOUNTED};
\node[com, right=of acc] (rec) {RECONCILED};
\node[st, below=9mm of rec] (ref) {REFUNDED};

\draw[arr] (created) -- (auth);
\draw[arr] (auth) -- (res);
\draw[arr] (res) -- (exec);
\draw[arr] (auth) to[out=-35,in=-145] node[below,pos=.5]{no reserve} (exec);
\draw[arr] (exec) -- (stream);
\draw[arr] (stream) -- (comp);
\draw[arr] (stream) -- (abort);
\draw[arr] (comp) -- (acc);
\draw[arr] (abort) -- (acc);
\draw[arr] (acc) -- (rec);
\draw[barr] (abort) to[out=-25,in=-155] node[below,pos=.5]{\color{red!70!black}M1: exit without commit} (ref);
\draw[arr] (res) to[out=60,in=120] node[above,pos=.5]{release} (rec);

\node[align=left, anchor=north west] at ($(created.south west)+(0,-14mm)$) {%
  \color{black!60}Shaded = value delivered\quad Darker = economic commitment};
\end{tikzpicture}%
}
\caption{Request lifecycle. Value exposure begins at the first delivered token
(\textsc{streaming}); economic commitment is final only at \textsc{reconciled}. The
interval between them is the exposure window that M1 exploits. B0 corrupts the
\textsc{authorized}$\rightarrow$\textsc{reserved} transition (non-atomic); M2 corrupts
the \emph{data} used at \textsc{accounted}.}
\label{fig:statemachine}
\end{figure}

\subsection{Formal model}

For a request $r$ let $U_{\text{true}}(r)$ be the usage actually produced,
$U_{\text{bill}}(r)$ the usage billed, and $P(\cdot)$ the (exact-decimal) pricing
function. Write $C(r)=P(U_{\text{true}}(r))$ for the true cost, $V(r)$ for the value
actually delivered to the client, $\mathrm{Res}(r)$ for any reservation, $D(r)$ for
committed debits, $F(r)$ for refunds, and
\[
N(r) \;=\; D(r)-F(r)
\]
for the \emph{net committed debit}. Under-payment and the integrity property are
\[
\mathrm{Leak}(r) \;=\; V(r)-N(r),
\qquad
\mathrm{Integrity}(r)\;:\;V(r)\le N(r).
\]
Because legitimate reconciliation exists, integrity is an inequality on the
\emph{net} debit; a refund is legitimate only when it releases the unconsumed portion,
\[
F(r)\;\le\;\mathrm{Res}(r)-V(r).
\]
At trial level the harness independently verifies conservation,
\[
\textstyle\sum_{r\in T} N(r) \;=\; B_{\text{initial}}-B_{\text{final}},
\]
and fails the build on any divergence.

\subsection{Mechanism localization}

\paragraph{B0 (synchronization).} If two concurrent requests read the same
authorization state $B$ and both commit a decrement derived from it, one decrement is
lost:
$\sum_i N(r_i) < \sum_i V(r_i)$. This requires $\ge 2$ concurrent transactions, which
is why B0 leaks nothing at concurrency~1.

\paragraph{M1 (commitment timing).} If $V(r)>0$ in \textsc{streaming} but the exit
path taken (\textsc{aborted}) does not reach \textsc{reconciled} with $N(r)\ge V(r)$,
then $\mathrm{Leak}(r)=V(r)>0$. The defect is in the transition graph, not in a shared
resource, hence it is concurrency-invariant.

\paragraph{M2 (usage authority).} If the billing basis is client-controlled,
$U_{\text{bill}}=T(U_{\text{true}})$ for an attacker transformation $T$, then
$\mathrm{Leak}(r)=P(U_{\text{true}})-P(T(U_{\text{true}}))>0$ whenever $P\circ T<P$.
Authority is a per-request data-provenance property, hence also
concurrency-invariant.

\begin{table}[t]\centering
\caption{Structural distinction between the three mechanisms.}
\label{tab:structural}
\resizebox{\linewidth}{!}{%
\begin{tabular}{llll}
\toprule
Mechanism & Defective element & Needs concurrency? & Closed by \\
\midrule
B0 & transition atomicity & \textbf{yes} ($\ge 2$) & atomic compare-and-decrement \\
M1 & transition graph (abort path) & no & abort-inclusive finalization \\
M2 & data authority at \textsc{accounted} & no & server-authoritative recount \\
\bottomrule
\end{tabular}%
}
\end{table}

This predicts the empirical result: B0 leakage is \emph{created} by concurrency, while
M1/M2 per-request leakage is invariant to it.
)
% Requires: tikz (with positioning, arrows.meta), amsmath, booktabs.

\section{Accounting State Model}
\label{sec:statemodel}

We model a metered inference request as a state machine so that each mechanism can be
localized to a specific transition, establishing that the taxonomy is
\emph{architectural} rather than a list of unrelated implementation bugs.

\begin{figure}[t]
\centering
\resizebox{\linewidth}{!}{%
\begin{tikzpicture}[
  node distance=8mm and 12mm,
  every node/.style={font=\scriptsize},
  st/.style={draw, rounded corners=2pt, minimum height=6mm, minimum width=17mm, align=center},
  val/.style={st, fill=black!6},
  com/.style={st, fill=black!12},
  bad/.style={st, draw=red!70!black, thick},
  ->/.default=,
  arr/.style={-{Latex[length=1.6mm]}, thin},
  barr/.style={-{Latex[length=1.6mm]}, thin, red!70!black, dashed},
]
\node[st] (created) {CREATED};
\node[st, right=of created] (auth) {AUTHORIZED};
\node[com, right=of auth] (res) {RESERVED};
\node[st, right=of res] (exec) {EXECUTING};
\node[val, right=of exec] (stream) {STREAMING};
\node[val, above right=4mm and 10mm of stream] (comp) {COMPLETED};
\node[bad, below right=4mm and 10mm of stream] (abort) {ABORTED};
\node[st, right=14mm of stream, xshift=26mm] (acc) {ACCOUNTED};
\node[com, right=of acc] (rec) {RECONCILED};
\node[st, below=9mm of rec] (ref) {REFUNDED};

\draw[arr] (created) -- (auth);
\draw[arr] (auth) -- (res);
\draw[arr] (res) -- (exec);
\draw[arr] (auth) to[out=-35,in=-145] node[below,pos=.5]{no reserve} (exec);
\draw[arr] (exec) -- (stream);
\draw[arr] (stream) -- (comp);
\draw[arr] (stream) -- (abort);
\draw[arr] (comp) -- (acc);
\draw[arr] (abort) -- (acc);
\draw[arr] (acc) -- (rec);
\draw[barr] (abort) to[out=-25,in=-155] node[below,pos=.5]{\color{red!70!black}M1: exit without commit} (ref);
\draw[arr] (res) to[out=60,in=120] node[above,pos=.5]{release} (rec);

\node[align=left, anchor=north west] at ($(created.south west)+(0,-14mm)$) {%
  \color{black!60}Shaded = value delivered\quad Darker = economic commitment};
\end{tikzpicture}%
}
\caption{Request lifecycle. Value exposure begins at the first delivered token
(\textsc{streaming}); economic commitment is final only at \textsc{reconciled}. The
interval between them is the exposure window that M1 exploits. B0 corrupts the
\textsc{authorized}$\rightarrow$\textsc{reserved} transition (non-atomic); M2 corrupts
the \emph{data} used at \textsc{accounted}.}
\label{fig:statemachine}
\end{figure}

\subsection{Formal model}

For a request $r$ let $U_{\text{true}}(r)$ be the usage actually produced,
$U_{\text{bill}}(r)$ the usage billed, and $P(\cdot)$ the (exact-decimal) pricing
function. Write $C(r)=P(U_{\text{true}}(r))$ for the true cost, $V(r)$ for the value
actually delivered to the client, $\mathrm{Res}(r)$ for any reservation, $D(r)$ for
committed debits, $F(r)$ for refunds, and
\[
N(r) \;=\; D(r)-F(r)
\]
for the \emph{net committed debit}. Under-payment and the integrity property are
\[
\mathrm{Leak}(r) \;=\; V(r)-N(r),
\qquad
\mathrm{Integrity}(r)\;:\;V(r)\le N(r).
\]
Because legitimate reconciliation exists, integrity is an inequality on the
\emph{net} debit; a refund is legitimate only when it releases the unconsumed portion,
\[
F(r)\;\le\;\mathrm{Res}(r)-V(r).
\]
At trial level the harness independently verifies conservation,
\[
\textstyle\sum_{r\in T} N(r) \;=\; B_{\text{initial}}-B_{\text{final}},
\]
and fails the build on any divergence.

\subsection{Mechanism localization}

\paragraph{B0 (synchronization).} If two concurrent requests read the same
authorization state $B$ and both commit a decrement derived from it, one decrement is
lost:
$\sum_i N(r_i) < \sum_i V(r_i)$. This requires $\ge 2$ concurrent transactions, which
is why B0 leaks nothing at concurrency~1.

\paragraph{M1 (commitment timing).} If $V(r)>0$ in \textsc{streaming} but the exit
path taken (\textsc{aborted}) does not reach \textsc{reconciled} with $N(r)\ge V(r)$,
then $\mathrm{Leak}(r)=V(r)>0$. The defect is in the transition graph, not in a shared
resource, hence it is concurrency-invariant.

\paragraph{M2 (usage authority).} If the billing basis is client-controlled,
$U_{\text{bill}}=T(U_{\text{true}})$ for an attacker transformation $T$, then
$\mathrm{Leak}(r)=P(U_{\text{true}})-P(T(U_{\text{true}}))>0$ whenever $P\circ T<P$.
Authority is a per-request data-provenance property, hence also
concurrency-invariant.

\begin{table}[t]\centering
\caption{Structural distinction between the three mechanisms.}
\label{tab:structural}
\resizebox{\linewidth}{!}{%
\begin{tabular}{llll}
\toprule
Mechanism & Defective element & Needs concurrency? & Closed by \\
\midrule
B0 & transition atomicity & \textbf{yes} ($\ge 2$) & atomic compare-and-decrement \\
M1 & transition graph (abort path) & no & abort-inclusive finalization \\
M2 & data authority at \textsc{accounted} & no & server-authoritative recount \\
\bottomrule
\end{tabular}%
}
\end{table}

This predicts the empirical result: B0 leakage is \emph{created} by concurrency, while
M1/M2 per-request leakage is invariant to it.
)
% Requires: tikz (with positioning, arrows.meta), amsmath, booktabs.

\section{Accounting State Model}
\label{sec:statemodel}

We model a metered inference request as a state machine so that each mechanism can be
localized to a specific transition, establishing that the taxonomy is
\emph{architectural} rather than a list of unrelated implementation bugs.

\begin{figure}[t]
\centering
\resizebox{\linewidth}{!}{%
\begin{tikzpicture}[
  node distance=8mm and 12mm,
  every node/.style={font=\scriptsize},
  st/.style={draw, rounded corners=2pt, minimum height=6mm, minimum width=17mm, align=center},
  val/.style={st, fill=black!6},
  com/.style={st, fill=black!12},
  bad/.style={st, draw=red!70!black, thick},
  ->/.default=,
  arr/.style={-{Latex[length=1.6mm]}, thin},
  barr/.style={-{Latex[length=1.6mm]}, thin, red!70!black, dashed},
]
\node[st] (created) {CREATED};
\node[st, right=of created] (auth) {AUTHORIZED};
\node[com, right=of auth] (res) {RESERVED};
\node[st, right=of res] (exec) {EXECUTING};
\node[val, right=of exec] (stream) {STREAMING};
\node[val, above right=4mm and 10mm of stream] (comp) {COMPLETED};
\node[bad, below right=4mm and 10mm of stream] (abort) {ABORTED};
\node[st, right=14mm of stream, xshift=26mm] (acc) {ACCOUNTED};
\node[com, right=of acc] (rec) {RECONCILED};
\node[st, below=9mm of rec] (ref) {REFUNDED};

\draw[arr] (created) -- (auth);
\draw[arr] (auth) -- (res);
\draw[arr] (res) -- (exec);
\draw[arr] (auth) to[out=-35,in=-145] node[below,pos=.5]{no reserve} (exec);
\draw[arr] (exec) -- (stream);
\draw[arr] (stream) -- (comp);
\draw[arr] (stream) -- (abort);
\draw[arr] (comp) -- (acc);
\draw[arr] (abort) -- (acc);
\draw[arr] (acc) -- (rec);
\draw[barr] (abort) to[out=-25,in=-155] node[below,pos=.5]{\color{red!70!black}M1: exit without commit} (ref);
\draw[arr] (res) to[out=60,in=120] node[above,pos=.5]{release} (rec);

\node[align=left, anchor=north west] at ($(created.south west)+(0,-14mm)$) {%
  \color{black!60}Shaded = value delivered\quad Darker = economic commitment};
\end{tikzpicture}%
}
\caption{Request lifecycle. Value exposure begins at the first delivered token
(\textsc{streaming}); economic commitment is final only at \textsc{reconciled}. The
interval between them is the exposure window that M1 exploits. B0 corrupts the
\textsc{authorized}$\rightarrow$\textsc{reserved} transition (non-atomic); M2 corrupts
the \emph{data} used at \textsc{accounted}.}
\label{fig:statemachine}
\end{figure}

\subsection{Formal model}

For a request $r$ let $U_{\text{true}}(r)$ be the usage actually produced,
$U_{\text{bill}}(r)$ the usage billed, and $P(\cdot)$ the (exact-decimal) pricing
function. Write $C(r)=P(U_{\text{true}}(r))$ for the true cost, $V(r)$ for the value
actually delivered to the client, $\mathrm{Res}(r)$ for any reservation, $D(r)$ for
committed debits, $F(r)$ for refunds, and
\[
N(r) \;=\; D(r)-F(r)
\]
for the \emph{net committed debit}. Under-payment and the integrity property are
\[
\mathrm{Leak}(r) \;=\; V(r)-N(r),
\qquad
\mathrm{Integrity}(r)\;:\;V(r)\le N(r).
\]
Because legitimate reconciliation exists, integrity is an inequality on the
\emph{net} debit; a refund is legitimate only when it releases the unconsumed portion,
\[
F(r)\;\le\;\mathrm{Res}(r)-V(r).
\]
At trial level the harness independently verifies conservation,
\[
\textstyle\sum_{r\in T} N(r) \;=\; B_{\text{initial}}-B_{\text{final}},
\]
and fails the build on any divergence.

\subsection{Mechanism localization}

\paragraph{B0 (synchronization).} If two concurrent requests read the same
authorization state $B$ and both commit a decrement derived from it, one decrement is
lost:
$\sum_i N(r_i) < \sum_i V(r_i)$. This requires $\ge 2$ concurrent transactions, which
is why B0 leaks nothing at concurrency~1.

\paragraph{M1 (commitment timing).} If $V(r)>0$ in \textsc{streaming} but the exit
path taken (\textsc{aborted}) does not reach \textsc{reconciled} with $N(r)\ge V(r)$,
then $\mathrm{Leak}(r)=V(r)>0$. The defect is in the transition graph, not in a shared
resource, hence it is concurrency-invariant.

\paragraph{M2 (usage authority).} If the billing basis is client-controlled,
$U_{\text{bill}}=T(U_{\text{true}})$ for an attacker transformation $T$, then
$\mathrm{Leak}(r)=P(U_{\text{true}})-P(T(U_{\text{true}}))>0$ whenever $P\circ T<P$.
Authority is a per-request data-provenance property, hence also
concurrency-invariant.

\begin{table}[t]\centering
\caption{Structural distinction between the three mechanisms.}
\label{tab:structural}
\resizebox{\linewidth}{!}{%
\begin{tabular}{llll}
\toprule
Mechanism & Defective element & Needs concurrency? & Closed by \\
\midrule
B0 & transition atomicity & \textbf{yes} ($\ge 2$) & atomic compare-and-decrement \\
M1 & transition graph (abort path) & no & abort-inclusive finalization \\
M2 & data authority at \textsc{accounted} & no & server-authoritative recount \\
\bottomrule
\end{tabular}%
}
\end{table}

This predicts the empirical result: B0 leakage is \emph{created} by concurrency, while
M1/M2 per-request leakage is invariant to it.
)
% Requires: tikz (with positioning, arrows.meta), amsmath, booktabs.

\section{Accounting State Model}
\label{sec:statemodel}

We model a metered inference request as a state machine so that each mechanism can be
localized to a specific transition, establishing that the taxonomy is
\emph{architectural} rather than a list of unrelated implementation bugs.

\begin{figure*}[t]
\centering
\resizebox{0.86\textwidth}{!}{%
\begin{tikzpicture}[
  node distance=8mm and 12mm,
  every node/.style={font=\scriptsize},
  st/.style={draw, rounded corners=2pt, minimum height=6mm, minimum width=17mm, align=center},
  val/.style={st, fill=black!6},
  com/.style={st, fill=black!12},
  bad/.style={st, draw=red!70!black, thick},
  ->/.default=,
  arr/.style={-{Latex[length=1.6mm]}, thin},
  barr/.style={-{Latex[length=1.6mm]}, thin, red!70!black, dashed},
]
\node[st] (created) {CREATED};
\node[st, right=of created] (auth) {AUTHORIZED};
\node[com, right=of auth] (res) {RESERVED};
\node[st, right=of res] (exec) {EXECUTING};
\node[val, right=of exec] (stream) {STREAMING};
\node[val, above right=4mm and 10mm of stream] (comp) {COMPLETED};
\node[bad, below right=4mm and 10mm of stream] (abort) {ABORTED};
\node[st, right=14mm of stream, xshift=26mm] (acc) {ACCOUNTED};
\node[com, right=of acc] (rec) {RECONCILED};
\node[st, below=9mm of rec] (ref) {REFUNDED};

\draw[arr] (created) -- (auth);
\draw[arr] (auth) -- (res);
\draw[arr] (res) -- (exec);
\draw[arr] (auth) to[out=-35,in=-145] node[below,pos=.5]{no reserve} (exec);
\draw[arr] (exec) -- (stream);
\draw[arr] (stream) -- (comp);
\draw[arr] (stream) -- (abort);
\draw[arr] (comp) -- (acc);
\draw[arr] (abort) -- (acc);
\draw[arr] (acc) -- (rec);
\draw[barr] (abort) to[out=-25,in=-155] node[below,pos=.5]{\color{red!70!black}M1: exit without commit} (ref);
\draw[arr] (res) to[out=60,in=120] node[above,pos=.5]{release} (rec);

\node[align=left, anchor=north west] at ($(created.south west)+(0,-14mm)$) {%
  \color{black!60}Shaded = value delivered\quad Darker = economic commitment};
\end{tikzpicture}%
}
\caption{Request lifecycle. Value exposure begins at the first delivered token
(\textsc{streaming}); economic commitment is final only at \textsc{reconciled}. The
interval between them is the exposure window that M1 exploits. B0 corrupts the
\textsc{authorized}$\rightarrow$\textsc{reserved} transition (non-atomic); M2 corrupts
the \emph{data} used at \textsc{accounted}.}
\label{fig:statemachine}
\end{figure*}

\subsection{Formal model}

For a request $r$ let $U_{\text{true}}(r)$ be the usage actually produced,
$U_{\text{bill}}(r)$ the usage billed, and $P(\cdot)$ the (exact-decimal) pricing
function. Write $C(r)=P(U_{\text{true}}(r))$ for the true cost, $V(r)$ for the value
actually delivered to the client, $\mathrm{Res}(r)$ for any reservation, $D(r)$ for
committed debits, $F(r)$ for refunds, and
\[
N(r) \;=\; D(r)-F(r)
\]
for the \emph{net committed debit}. Under-payment and the integrity property are
\[
\mathrm{Leak}(r) \;=\; V(r)-N(r),
\qquad
\mathrm{Integrity}(r)\;:\;V(r)\le N(r).
\]
Because legitimate reconciliation exists, integrity is an inequality on the
\emph{net} debit; a refund is legitimate only when it releases the unconsumed portion,
\[
F(r)\;\le\;\mathrm{Res}(r)-V(r).
\]
At trial level the harness independently verifies conservation,
\[
\textstyle\sum_{r\in T} N(r) \;=\; B_{\text{initial}}-B_{\text{final}},
\]
and fails the build on any divergence.

\subsection{Mechanism localization}

\paragraph{B0 (synchronization).} In the baseline implementation, if two concurrent
requests read the same authorization state $B$ and both commit a decrement derived from
it, one decrement is lost: $\sum_i N(r_i) < \sum_i V(r_i)$. That particular
manifestation requires $\ge 2$ overlapping transactions, which is why the baseline leaks
nothing at concurrency~1. The general failure condition is broader than its trigger: it
holds whenever authorization observes economic state that does not yet include the
commitments produced by previously authorized service. Overlapping requests are one way
to create that separation; \S\ref{sec:b0revised} shows that deferred settlement creates
the same condition under strictly sequential arrivals. The condition is therefore that
authorization be atomically coupled to economic commitment on the path that authorizes
service, and the baseline trigger must not be mistaken for the condition itself.

\paragraph{M1 (commitment timing).} If $V(r)>0$ in \textsc{streaming} but the exit
path taken (\textsc{aborted}) does not reach \textsc{reconciled} with $N(r)\ge V(r)$,
then $\mathrm{Leak}(r)=V(r)>0$. The defect is in the transition graph, not in a shared
resource, so the model gives no reason for concurrency to be causal here; under fixed
request volume we measured identical per-request leakage across concurrency 1--100 in
the tested architecture (\S\ref{sec:m1}). We report that as a measurement, not as a
property of all deployments.

\paragraph{M2 (usage authority).} If the billing basis is client-controlled,
$U_{\text{bill}}=T(U_{\text{true}})$ for an attacker transformation $T$, then
$\mathrm{Leak}(r)=P(U_{\text{true}})-P(T(U_{\text{true}}))>0$ whenever $P\circ T<P$.
Authority is a per-request data-provenance property: under the modelled per-request
billing function the billed amount reads no shared mutable state, so independence from
concurrency is analytic, and the sweep serves as an implementation check rather than as
its evidence.

\begin{table}[t]\centering
\caption{Structural distinction between the three mechanisms. ``Baseline trigger'' is
what exercises the defect in the implementation we measured; it is not the security
condition. For B0 in particular, overlapping requests are one trigger and deferred
settlement is another (\S\ref{sec:b0revised}), so the condition is stated in terms of
authorization/commitment coupling rather than concurrency.}
\label{tab:structural}
\resizebox{\linewidth}{!}{%
\begin{tabular}{llll}
\toprule
Mechanism & Defective element & Baseline trigger & Closed by \\
\midrule
B0 & transition atomicity & overlapping requests ($\ge 2$), or deferred settlement & authorization atomically coupled to commitment on the authorizing path \\
M1 & transition graph (abort path) & a single abort; no concurrency required & abort-inclusive finalization \\
M2 & data authority at \textsc{accounted} & a single request; no concurrency required & server-authoritative recount \\
\bottomrule
\end{tabular}%
}
\end{table}

This predicts the empirical result: the baseline B0 manifestation is \emph{created} by
concurrency, while M1/M2 per-request leakage is invariant to it. \S\ref{sec:b0revised}
then shows that the concurrency dependence belonged to that baseline trigger rather than
to the underlying condition.
